\documentclass{article}

\usepackage{amsmath,amssymb}
\usepackage{xurl}
\usepackage[hidelinks]{hyperref}
\setlength{\emergencystretch}{2em}

% Shared commands for notes in docs/paper-gaps/.

\DeclareUnicodeCharacter{2080}{\ifmmode{}_{0}\else\textsubscript{0}\fi}
\DeclareUnicodeCharacter{2081}{\ifmmode{}_{1}\else\textsubscript{1}\fi}
\DeclareUnicodeCharacter{2082}{\ifmmode{}_{2}\else\textsubscript{2}\fi}
\DeclareUnicodeCharacter{2099}{\ifmmode{}_{n}\else\textsubscript{n}\fi}

\newcommand{\Lean}{\textsc{Lean}}
\ExplSyntaxOn
\NewDocumentCommand{\leanid}{m}{
  \ifmmode
    \textsf{\detokenize{#1}}
  \else
    \group_begin:
    \urlstyle{sf}
    \tl_set:Nn \l_tmpa_tl {#1}
    \tl_replace_all:Nnn \l_tmpa_tl {\_} {_}
    \exp_args:NV \path \l_tmpa_tl
    \group_end:
  \fi
}
\ExplSyntaxOff
\providecommand{\tr}{\operatorname{tr}}
\newcommand{\tnleanIssueBase}{https://github.com/LionSR/TNLean/issues/}
\newcommand{\tnleanPrBase}{https://github.com/LionSR/TNLean/pull/}
\newcommand{\tnleanDocBase}{https://sirui-lu.com/TNLean/docs/}
\newcommand{\ghissue}[1]{\href{\tnleanIssueBase#1}{GitHub issue \##1}}
\newcommand{\ghpr}[1]{\href{\tnleanPrBase#1}{PR \##1}}
\newcommand{\doclink}[1]{\href{\tnleanDocBase#1.html}{\path{#1}}}

% Dirac notation and the identity operator, as in blueprint/src/macros/common.tex.
% \providecommand keeps note-local definitions authoritative.
\usepackage{dsfont}
\providecommand{\ket}[1]{|#1\rangle}
\providecommand{\bra}[1]{\langle #1|}
\providecommand{\braket}[2]{\langle #1 | #2 \rangle}
\providecommand{\ketbra}[2]{|#1\rangle\!\langle #2|}
\providecommand{\Id}{\mathds{1}}
\providecommand{\one}{\mathds{1}}
\providecommand{\C}{\mathbb{C}}
\providecommand{\MN}[1]{M_{#1}(\C)}
\providecommand{\MD}{M_D(\C)}

% External referencing for paper sources.  The build helper writes sanitized
% auxiliary files under build/paper-aux/ and excludes duplicated source labels.
\usepackage{xr}

% Import a generated paper auxiliary file with a caller-supplied label prefix.
% The candidate paths support builds from docs/paper-gaps/, the repository root,
% and build/paper-aux/ itself.
\newcommand{\importpaperaux}[2]{%
  \IfFileExists{../../build/paper-aux/#2.aux}{%
    \externaldocument[#1]{../../build/paper-aux/#2}%
  }{%
    \IfFileExists{build/paper-aux/#2.aux}{%
      \externaldocument[#1]{build/paper-aux/#2}%
    }{%
      \IfFileExists{#2.aux}{%
        \externaldocument[#1]{#2}%
      }{}%
    }%
  }%
}

% General external reference macros
\newcommand{\paperref}[2]{\ref{#1-#2}}
\newcommand{\papereqref}[2]{(\ref{#1-#2})}

% CPSV16 source (1606.00608 / MPDO-22-12-17-2.tex) external document and shortcuts
\importpaperaux{cpsv16-}{1606.00608/MPDO-22-12-17-2-tnlean}
\newcommand{\cpsvref}[1]{\paperref{cpsv16}{#1}}
\newcommand{\cpsveqref}[1]{\papereqref{cpsv16}{#1}}

% Verdict marker: \gapnote{<kind>}{<status>}. Kind is the policy
% classification (clarification, local-correction, scope-restriction,
% unfaithful, false-source, open-gap); status is open, wip (elimination
% underway), resolved, or historical. Machine-read by the site index;
% prints a verdict line.
\newcommand{\gapnote}[2]{\par\noindent\textbf{Verdict:} #1 (#2).\par}


\title{PEPS Gauge Uniqueness and Balanced Edge Scalars}
\author{}
\date{2026-04-30}

\begin{document}
\maketitle
\gapnote{false-source}{resolved}

\section{The assertion}

The uniqueness clause of the Fundamental Theorem for injective PEPS on a
simple graph appears in Theorem~2 and Section~3 of
Ref.~\cite{Molnar2018NormalPEPS}. The proof first blocks the graph around a
chosen edge into a three-site injective MPS. The one-dimensional fundamental
theorem then assigns a gauge to that edge. Repeating this construction over all
edges gives a family of edge gauges, and a local two-tensor comparison gives
scalar proportionalities at the vertices.\footnote{In the repository source,
the edge-blocking step is at
\path{Papers/1804.04964/paper_normal.tex:981--1009}; the assignment of edge
gauges is at \path{Papers/1804.04964/paper_normal.tex:1037--1065}; the local
scalar proportionality argument is at
\path{Papers/1804.04964/paper_normal.tex:1066--1092} and
\path{Papers/1804.04964/paper_normal.tex:1204--1207}.}

The theorem applies to PEPS on graphs without double edges or self-loops. It
states that two injective PEPS generate the same state if and only if their
tensors are related by local gauges. The source summarizes the uniqueness
claim by saying that the local gauges are ``unique up to a multiplicative
constant'' \cite{Molnar2018NormalPEPS}.\footnote{The quoted phrase occurs at
\path{Papers/1804.04964/paper_normal.tex:1207--1210}. For traceability, this
note corresponds to \ghissue{1045}.}
If this phrase means uniqueness up to one common scalar multiplying every edge
gauge, the claim is false on a general graph.

\section{The corrected quotient}

Let $G=(V,E)$ be a finite simple graph with a fixed orientation on each edge.
For every $e\in E$, let $X_e$ be an invertible matrix on the bond space of
$e$. At the tail of $e$, the endpoint action is $X_e$; at the head, it is
$X_e^{-1}$.

Let $c=(c_e)_{e\in E}$ be a family of nonzero edge scalars. Define its oriented
endpoint value by
\begin{align}
    c_{v,e}
        &=
        \begin{cases}
            c_e, & \text{if $v$ is the tail of $e$},\\
            c_e^{-1}, & \text{if $v$ is the head of $e$}.
        \end{cases}
    \label{eq:peps_gauge_edge_scalars_oriented_endpoint_scalar}
\end{align}
The family $c$ is vertex-balanced when
\begin{align}
    \prod_{e\ni v}c_{v,e}
        &= 1
        \qquad\text{for every }v\in V.
    \label{eq:peps_gauge_edge_scalars_vertex_balance}
\end{align}

For gauge families $X$ and $Y$, write $X_{v,e}$ and $Y_{v,e}$ for their
oriented endpoint actions. The two families are equivalent modulo balanced
edge scalars if there is a vertex-balanced family $c$ such that
\begin{align}
    X_{v,e}
        &= c_{v,e}Y_{v,e}
        \qquad\text{for every incident pair }(v,e).
    \label{eq:peps_gauge_edge_scalars_quotient_relation}
\end{align}
This is the graph-correct quotient. Multiplying the incident endpoint gauges at
$v$ by $c_{v,e}$ multiplies the gauged local tensor by
\begin{align}
    \prod_{e\ni v}c_{v,e}.
    \label{eq:peps_gauge_edge_scalars_local_tensor_multiplier}
\end{align}
For a vertex-balanced family, this multiplier is one by
Eq.~\ref{eq:peps_gauge_edge_scalars_vertex_balance}, so every gauged local
tensor is unchanged. Gauge uniqueness can therefore hold at most modulo
balanced edge scalars.

The formal statement uses this quotient as the corrected uniqueness relation,
and the relation is an equivalence relation on gauge families.\footnote{The
formal definitions are \leanid{TNLean.PEPS.edgeScalarAt},
\leanid{TNLean.PEPS.IsVertexBalanced}, and
\leanid{TNLean.PEPS.GaugeEquivModEdgeScalars}. The quotient relation is
\leanid{TNLean.PEPS.GaugeEquivModEdgeScalars.setoid} in
\path{TNLean/PEPS/FundamentalTheorem.lean}. The corresponding blueprint entries
are in \path{blueprint/src/chapter/ch24_peps_ft.tex}.}

\section{The triangle witness}

The obstruction already occurs on a connected triangle. Let the vertices be
ordered by $1<2<3$, orient each edge from its smaller endpoint to its larger
endpoint, and set every bond dimension to one. An edge gauge is then a nonzero
complex number.

Take any nonzero injective bond-dimension-one PEPS tensor $A$. Compare the
identity gauge family with the scalar gauge family
\begin{align}
    c_{12}
        &= t,
    \label{eq:peps_gauge_edge_scalars_triangle_edge_twelve}\\
    c_{23}
        &= t,
    \label{eq:peps_gauge_edge_scalars_triangle_edge_twentythree}\\
    c_{13}
        &= t^{-1},
    \label{eq:peps_gauge_edge_scalars_triangle_edge_thirteen}\\
    t
        &\in \C^\times,
        \qquad t^2\neq 1.
    \label{eq:peps_gauge_edge_scalars_triangle_parameter}
\end{align}
The oriented products at the three vertices are
\begin{align}
    c_{12}c_{13}
        &= tt^{-1}
         = 1,
    \label{eq:peps_gauge_edge_scalars_triangle_balance_vertex_one}\\
    c_{12}^{-1}c_{23}
        &= t^{-1}t
         = 1,
    \label{eq:peps_gauge_edge_scalars_triangle_balance_vertex_two}\\
    c_{13}^{-1}c_{23}^{-1}
        &= tt^{-1}
         = 1.
    \label{eq:peps_gauge_edge_scalars_triangle_balance_vertex_three}
\end{align}
Thus the family is vertex-balanced. Applying it to $A$ leaves every local
tensor unchanged, so it gives a second gauge family relating $A$ to itself.

The two gauge families do not differ by one global scalar. If a scalar
$s\in\C^\times$ satisfied
\begin{align}
    c_{12}
        &= c_{13}
         = c_{23}
         = s,
    \label{eq:peps_gauge_edge_scalars_hypothetical_global_scalar}
\end{align}
then $c_{12}=c_{13}$ would imply
\begin{align}
    t
        &= t^{-1},
    \label{eq:peps_gauge_edge_scalars_triangle_inverse_equality}\\
    t^2
        &= 1,
    \label{eq:peps_gauge_edge_scalars_triangle_parameter_contradiction}
\end{align}
contrary to Eq.~\ref{eq:peps_gauge_edge_scalars_triangle_parameter}. Therefore
uniqueness modulo one global scalar is false even for a connected simple graph
with one-dimensional bonds. The balanced edge-scalar quotient identifies the
two families exactly.\footnote{This counterexample is recorded in
\path{docs/counterexamples.md:127--135}. It originated in \ghissue{762}; the
formal statement was corrected in \ghpr{790}.}

\section{Formal statement}

The discrepancy changes the statement, not merely its proof. The triangle
example satisfies the graph-theoretic hypotheses of the uniqueness clause in
Ref.~\cite{Molnar2018NormalPEPS} and gives two genuine gauge families which
differ by a balanced edge-scalar family but not by one global scalar. The
source phrase must therefore be read as implicitly referring to the balanced
quotient or corrected to state uniqueness modulo balanced edge scalars.

The corrected quotient is proved under the standing normal-PEPS convention
that every bond space is nonzero. Equality of the gauged local tensors at a
vertex first gives equality between products of incident edge-gauge entries for
every pair of virtual configurations. Linear independence then promotes this
coefficient equality to the required product identity.

Each oriented endpoint gauge is invertible. Because every bond space is
nonzero, its index type is nonempty. The many-leg scalar extraction therefore
gives nonzero endpoint proportionality constants $c_{v,e}$ satisfying
\begin{align}
    X_{v,e}
        &= c_{v,e}Y_{v,e}
        \qquad\text{for every incident pair }(v,e),
    \label{eq:peps_gauge_edge_scalars_formal_endpoint_proportionality}\\
    \prod_{e\ni v}c_{v,e}
        &= 1
        \qquad\text{for every }v\in V.
    \label{eq:peps_gauge_edge_scalars_formal_vertex_product}
\end{align}
The lower-endpoint constant of each edge defines one edge scalar. The inverse
relation transports it to the upper endpoint. Invertibility makes each
constant unique, so
Eq.~\ref{eq:peps_gauge_edge_scalars_formal_vertex_product} becomes the
oriented vertex-balance condition in
Eq.~\ref{eq:peps_gauge_edge_scalars_vertex_balance}.\footnote{For traceability,
the declaration is
\leanid{TNLean.PEPS.gauge_unique_mod_edge_scalars}; the many-leg extraction is
\leanid{TNLean.PEPS.piProduct_forms_scalar} in
\path{TNLean/PEPS/TensorFactorScalar.lean}. This point is recorded in
\ghissue{842}; \ghpr{879} reduced the formal statement to the scalar-ratio and
global-reconciliation step, which has since been completed.}

\subsection{Nonzero bond spaces}

The hypothesis
\begin{align}
    \dim_e
        &> 0
        \qquad\text{for every }e\in E
    \label{eq:peps_gauge_edge_scalars_positive_bond_dimension}
\end{align}
is necessary. If a bond space is empty, the local virtual-configuration types
at both endpoints are empty, so the corresponding linear-independence
hypotheses are vacuous. A positive-dimensional bond adjacent to an endpoint
that also meets such an empty bond is then unconstrained by those hypotheses.

One may consequently choose two gauge families that agree on every other bond
but differ on that positive-dimensional bond by no scalar. Both families
satisfy the relevant hypotheses vacuously, while the balanced edge-scalar
conclusion fails. The nonzero-bond convention is precisely the normal-PEPS
assumption used in the existence direction, where it appears as
\leanid{hposA} in \leanid{TNLean.PEPS.fundamentalTheorem_PEPS}. The uniqueness
clause requires the same convention.

\section{Verdict}

The source's global-scalar uniqueness interpretation is false as printed. The
connected triangle provides a counterexample with one-dimensional bond spaces.
The corrected statement is uniqueness modulo vertex-balanced edge scalars
under the nonzero-bond-space convention. This quotient preserves every gauged
local tensor and is now formalized and proved.

\bibliographystyle{alpha}
\bibliography{references}

\end{document}
